

Reference blink events for both corpora were established by manual annotation. Two raters, blinded to the study objectives and to the outputs of the blink-detection algorithms evaluated in this study, annotated the recordings, and each blink was marked as an interval delimited by its onset and offset. The first rater identified the blinks in each recording and marked their onsets and offsets. The second rater then reviewed every recording independently, checking the onset and offset of each marked interval and searching the recording for blinks that had been missed or events that had been marked incorrectly. Discrepant annotations were reviewed jointly and resolved by consensus, and the consensus labels served as the reference for all evaluations. Because the second rater reviewed the first rater's annotations rather than producing a separate annotation from scratch, no inter-rater agreement statistic was computed. A recording entered the analysis only after its annotation had been marked as complete.

\paragraph{Internal corpus.}
For the Internal corpus, a face video recorded in synchrony with the EEG served as the primary behavioural reference. The video was captured by a USB webcam at 30 frames/s and a resolution of $1280\times720$ pixels.
%The acquisition software (Net Station 5.1.2) limited each video file to 30~min and wrote, for each file, a synchronization record mapping the video onto EEG sample time. Each session of the Internal corpus therefore corresponds to one such video segment.
Annotation relied on three time-aligned sources, namely the video itself, a vertical electrooculographic (EOG) derivation, and an eye aspect ratio (EAR) trace computed from facial landmarks in each video frame~\citep{dewi2022adjusting}. The EAR quantifies eye opening, falling as the eyelids close and recovering as they reopen. An event was accepted as a blink only when a closing and reopening of the eyelids was visible in the video. Blink onset was placed where the eyelids began to close and the EAR began to fall from its open-eye level, and blink offset where the eyelids had returned to the open position and the EAR had recovered. Each boundary was then checked against the characteristic high-amplitude blink transient in the vertical EOG~\citep{croft2000removal}. The video and EAR thus determined whether an event was a blink and where it began and ended, whereas the vertical EOG provided an electrophysiological check on that timing. All marked events, whether the eyelid closure was complete or partial, were treated as reference blinks in the evaluation.

\paragraph{Cao2018 corpus.}
No video recordings or separate EOG channels are distributed with the Cao2018 corpus. Its event markers describe only the driving task, namely lane-departure onsets and the onsets and offsets of the driver's corrective responses, and no blink annotations are provided~\citep{cao2019multichannel}. Blinks in this corpus were therefore annotated from the EEG alone. The raters inspected each recording with Fp1 and Fp2 displayed by default and the remaining channels available on demand. They marked as blinks the transient, high-amplitude deflections ~\citep{iwasaki2005effects,croft2000removal,kleifges2017blinker}. Blink onset was placed where this deflection departed from the preceding baseline, and blink offset where it returned to baseline.
%Because these labels rest on EEG morphology alone, they are expected to be less certain than the video-supported Internal labels, particularly for small-amplitude blinks and for ocular events with a similar waveform. Four Cao2018 sessions that were flagged during annotation as having unresolved issues were excluded, leaving 58 of the 62 sessions for analysis.
